% Insert after Section 2. No document environment is included.
% Required preamble package: \usepackage{hyperref}
% Current implementation checkpoint: main, dbbe536 (30 October 2026).
% Earlier Thunder Client scenarios/screenshots: supplied Section 3 document,
% which identifies feat/features, 8da171a as the tested backend checkpoint.
% Do not present these manual scenarios as new tests on dbbe536.
% Current sources: docs/MVP.md; docs/validation/MVP-checks-2026-10-30.md;
% frontend/docs/interim-student-checks.md; current backend/frontend source.
% Retain selected screenshots from the supplied document beside the relevant
% scenarios, with their original tested checkpoint identified in the captions.

\section{Implemented Core Features and Integration Evidence}
\label{sec:core}

The interim implementation checkpoint is \texttt{main}, commit
\texttt{dbbe536}, dated 30 October 2026. It integrates authentication,
course learning, submissions, grading, instructor review, recommendations
and management interfaces. The Thunder Client scenarios below document
earlier backend checks at commit \texttt{8da171a}; they are distinguished
from the later build and automated validation evidence. An implemented
feature is not treated as fully verified solely because its code or
interface exists.

\subsection{Mandatory Functional Coverage}

\subsubsection{Accounts and Authorisation}
Registration, login, logout, refresh-token rotation, profile retrieval and
editing, role checks and ownership protection are implemented. Management
functions support account roles and status changes with audit records.
Automated checks cover selected permission and session behaviours; earlier
API demonstrations also exercised authentication failures and access denial.
Password recovery remains outside the implemented authentication flow, and
complete browser acceptance of account-management operations is pending.

\subsubsection{Courses and Content}
Course creation, editing, publication and hiding, chapters, text lessons and
Writing activities are implemented. The learner catalogue provides search,
filtering and pagination over the fetched course list. Lesson rendering
supports text and external links, while the authoring workflow primarily
creates text content. End-to-end authoring and delivery of a second material
type, and complete presentation/enforcement of course prerequisites, still
require verification. Skill prerequisites used by recommendations do not
alone satisfy course-prerequisite requirements.

\subsubsection{Enrolment and Learning}
Learners can enrol, leave courses, open permitted lessons and persist lesson
completion. The interface displays saved progress, and configured course
completion checks are implemented. Earlier API checks demonstrated that
leaving a course prevents further lesson access. Selected learner browser
checks and frontend regressions provide additional evidence. Resuming the
last study position and a complete in-progress/completed-course overview
remain integration or acceptance tasks.

\subsubsection{Assignments and Assessment}
Versioned assignments, time limits, deadlines, draft locking, immutable
submission revisions, objective grading, instructor feedback, publication
and attempt history are implemented. The Reading example uses answer-key
grading; Writing supports teacher review and revision. Integration tests
cover repeated submissions, stale updates and publication behaviour.
Provider transport tests use mocks, so live AI assessment and scoring
accuracy are not established. File/audio handling is outside the verified
demonstration paths.

\subsubsection{Interaction and Notifications}
Assessment publication creates notifications; retrieval, read-state updates
and a notification page are implemented. Appeals preserve learner reasons
and instructor resolutions. These functions do not replace the required
discussion, Q\&A or comment channel, or inappropriate-content reporting,
which remain outstanding. Notification coverage for new content and
upcoming deadlines also requires completion or verification.

\subsubsection{Administration and Reporting}
The implementation includes account/role/status management, instructor
application decisions, course-instructor assignment, skill administration
and audit recording. Role-scoped operational reports summarise grading,
review and provider activity. These reports do not establish complete
enrolment, completion-rate and learning-outcome analytics. A full
administrator browser scenario, including management changes and their
audit evidence, remains to be demonstrated.

\subsection{Integrated Demonstration Scenarios}

\subsubsection{Learner Scenario}
In the earlier Thunder Client demonstration, a test learner signed in,
enrolled in a published course, accessed a lesson and recorded completion.
A Reading attempt used a 300-second limit and version-controlled draft
updates. Submission created an immutable revision and queued objective
grading. Correct answers produced \textbf{2/2 RAW marks}, not an IELTS band.
Repeating the submission returned the existing revision; another learner's
request for that submission returned HTTP~403. The supplied screenshots
show the RAW result and rejected access. Integration tests additionally
cover repeated submissions and stale draft updates.

% Suggested retained caption: Earlier backend demonstration at 8da171a:
% answer-key assessment returns RAW 2/2; unauthorised access returns HTTP 403.

\subsubsection{Instructor Scenario}
The earlier Writing demonstration used synthetic text, a test rubric and
\texttt{MANUAL} result release. Without an AI configuration, submission
entered \texttt{AWAITING\_REVIEW}. The assigned instructor claimed the case
and recorded four criterion scores, feedback and a skill observation.
Before publication, the learner received an empty assessment list. After
publication, the teacher assessment became visible. The before/after
screenshots support this release-policy behaviour. Integration tests also
check that unpublished results do not expose accepted skill evidence.
Synthetic inputs and teacher-entered scores demonstrate workflow behaviour,
not AI accuracy.

% Suggested retained caption: Earlier MANUAL-release demonstration at
% 8da171a: assessment visibility before and after instructor publication.

\subsubsection{Appeal Handling and Administrator Verification}
The earlier API test record reports a learner appeal, rejection of a
duplicate open appeal, and instructor resolution with a persisted reason.
In the tested rejection path, the original score remained unchanged.
These actions demonstrate learner--instructor integration rather than an
administrator workflow. At the later checkpoint, local API smoke checks
exercise all three roles, including rejection of learner access to
operational reports. Administrator management interfaces are implemented,
but a complete browser demonstration of an account/course management action,
its audit record and reporting remains an acceptance task.

\subsection{Integration Status and Known Functional Gaps}

The current implementation connects authenticated users to enrolments,
learning progress, pinned assignment versions, submission revisions,
assessments and review decisions in PostgreSQL. Only accepted and released
evidence contributes to the official skill profile; recommendations apply
eligibility and prerequisite checks to that profile. Frontend pages are
connected to these APIs, although full three-role browser acceptance remains
pending.

The
\href{https://github.com/huyho44/IELTS-Evaluator-Pro/blob/dbbe536/docs/validation/MVP-checks-2026-10-30.md}{30 October validation record}
reports 19 passing isolated database integration tests, 11 local API smoke
checks, five passing frontend regression scripts and successful production
builds. Migration reproducibility was checked by applying three migrations
and the seed to a fresh isolated test database. The earlier backend
implementation is traceable through
\href{https://github.com/huyho44/IELTS-Evaluator-Pro/commit/8da171afdce3bfe58bec1d10ded401ce745ad9a3}{commit 8da171a}.

Remaining work includes the mandatory communication/reporting gaps,
password recovery, complete content and study-flow acceptance, public
deployment, and independent evaluation of assessment and recommendations.
Section~\ref{sec:plan} sets out proposed responsibilities and completion
milestones, prioritising mandatory requirements over optional extensions.
